\newpage
\section{Kết luận}

Trong đồ án này, nhóm chúng em đã tiến hành nghiên cứu, phân tích và thực nghiệm chi tiết về chủ đề Tổng hợp khuôn mặt theo định hướng từ giáo trình môn Nhận dạng. Cụ thể, nhóm đã trình bày một cái nhìn tổng quan về quá trình phát triển của lĩnh vực này, từ các phương pháp truyền thống đến các mô hình học sâu hiện đại. Trọng tâm của đồ án là việc tìm hiểu sâu và tái hiện lại hệ thống tiên tiến SynFER. Dù gặp nhiều hạn chế về mặt tài nguyên máy tính cũng như việc không có sẵn trọng số huấn luyện từ nhóm tác giả gốc, nhóm đã nỗ lực tự huấn luyện mô hình từ đầu và thành công xây dựng một quy trình sinh ảnh đầy đủ. Kết quả thực nghiệm cho thấy mặc dù chưa đạt được độ hội tụ tối đa về mặt chất lượng hình ảnh so với công bố ban đầu, mô hình vẫn thể hiện khả năng kiểm soát biểu cảm linh hoạt, độ chính xác nhận dạng cao, đồng thời làm rõ được những điểm mạnh và điểm nghẽn kỹ thuật khi áp dụng phương pháp này vào thực tế. Việc xây dựng thành công ứng dụng demo minh họa cũng chứng tỏ tính ứng dụng to lớn của giải pháp.

